% The Preamble is where you define the document type and load packages.

\documentclass{report}

% Load packages that add new features
\usepackage[utf8]{inputenc} % Allows you to use characters like á, ç, ñ directly
\usepackage{amsmath}        % For advanced math environments
\usepackage{graphicx}       % To include images

% Document Metadata
\title{Diagnóstico de Falhas em Máquinas Rotativas por Meio de Simulações Numéricas e Análise de Espectros de Ordem Superior}

\author{Cristofer Antoni Souza Costa}

\date{\today}

% --- This is the Document Body ---
% The content of your document goes between \begin{document} and \end{document}.

\begin{document}

\maketitle % This command displays the title, author, and date

\section{Introduction} % Creates a numbered section heading

% =============================================================================
% CAPÍTULO 1: INTRODUÇÃO
% =============================================================================

\chapter{Introdução}

\section{Contextualização}

As máquinas rotativas constituem elementos fundamentais em diversos setores industriais, desde a geração de energia até processos de manufatura. Sua operação confiável é essencial para a eficiência e segurança dos processos industriais, sendo que falhas não detectadas podem resultar em paradas não planejadas, perdas de produção e, em casos extremos, acidentes com consequências graves.

O monitoramento e diagnóstico de falhas em máquinas rotativas tem evoluído significativamente nas últimas décadas, impulsionado pelo desenvolvimento de técnicas avançadas de análise de sinais e pela crescente demanda por sistemas de manutenção preditiva. Neste contexto, a análise de vibrações emerge como uma das principais ferramentas para o diagnóstico de falhas, oferecendo informações valiosas sobre o estado operacional dos equipamentos.

\section{Justificativa}

A análise de espectros de ordem superior (Higher-Order Spectra - HOS), incluindo bi-espectro e tri-espectro, oferece uma abordagem promissora para identificar assinaturas não lineares associadas a diferentes tipos de falhas. No entanto, a aplicação prática dessas técnicas muitas vezes enfrenta desafios devido à complexidade dos fenômenos envolvidos e à necessidade de dados experimentais extensivos.

Este trabalho propõe o desenvolvimento de uma metodologia baseada em simulações numéricas para o diagnóstico de falhas em máquinas rotativas, utilizando a análise de HOS. A abordagem por simulação oferece vantagens significativas, incluindo controle preciso dos parâmetros do sistema, capacidade de gerar grandes volumes de dados e possibilidade de estudar cenários que seriam difíceis ou perigosos de reproduzir experimentalmente.

\section{Problema de Pesquisa}

O problema central desta pesquisa consiste em desenvolver uma metodologia eficaz para o diagnóstico de falhas em máquinas rotativas utilizando simulações numéricas e análise de espectros de ordem superior, com o objetivo de identificar padrões característicos que permitam distinguir entre condições normais de operação e diferentes tipos de falhas.

\section{Objetivos}

\subsection{Objetivo Geral}

Desenvolver uma metodologia baseada em simulações numéricas para o diagnóstico de falhas em máquinas rotativas utilizando análise de espectros de ordem superior.

\subsection{Objetivos Específicos}

\begin{enumerate}
    \item Desenvolver modelos de elementos finitos utilizando a ferramenta ROSS para simular o comportamento dinâmico de eixos rotativos sob condições normais e com falhas específicas, como trincas e desalinhamentos.
    
    \item Aplicar técnicas de análise de espectros de ordem superior (bi-espectro e tri-espectro) aos dados simulados para identificar padrões característicos que diferenciem as condições normais das falhas modeladas.
    
    \item Validar os resultados obtidos através de comparação com dados disponíveis na literatura científica.
    
    \item Avaliar a eficácia da metodologia proposta para diferentes tipos de falhas e condições operacionais.
\end{enumerate}

\section{Organização do Trabalho}

Esta dissertação está organizada em cinco capítulos, além dos elementos pré e pós-textuais:

\begin{itemize}
    \item \textbf{Capítulo 1 - Introdução:} Apresenta a contextualização, justificativa, problema de pesquisa, objetivos e organização do trabalho.
    
    \item \textbf{Capítulo 2 - Revisão de Literatura:} Aborda os fundamentos teóricos relacionados à dinâmica de rotores, análise de vibrações, espectros de ordem superior e diagnóstico de falhas.
    
    \item \textbf{Capítulo 3 - Metodologia:} Descreve detalhadamente a metodologia proposta, incluindo a modelagem numérica, simulações e análise de dados.
    
    \item \textbf{Capítulo 4 - Resultados e Discussão:} Apresenta e discute os principais resultados obtidos, incluindo a análise dos espectros de ordem superior e a identificação de padrões característicos.
    
    \item \textbf{Capítulo 5 - Conclusões:} Sintetiza os principais achados, contribuições e perspectivas futuras do trabalho.
\end{itemize}


% =============================================================================
% CAPÍTULO 2: REVISÃO DE LITERATURA
% =============================================================================

\chapter{Revisão de Literatura}

\section{Dinâmica de Rotores}

A dinâmica de rotores é uma área fundamental da engenharia mecânica que estuda o comportamento dinâmico de sistemas rotativos. Segundo Friswell et al. \cite{friswell2010}, os sistemas rotativos apresentam características dinâmicas únicas que os distinguem de outros sistemas mecânicos, incluindo efeitos giroscópicos, forças centrífugas e acoplamentos entre movimentos translacionais e rotacionais.

\subsection{Fundamentos Teóricos}

A modelagem matemática de sistemas rotativos baseia-se nas equações de movimento de Lagrange ou Hamilton, considerando as forças inerciais, de amortecimento e de rigidez. Para um rotor simples, a equação de movimento pode ser expressa como:

\begin{equation}
\mathbf{M}\ddot{\mathbf{q}} + (\mathbf{C} + \mathbf{G})\dot{\mathbf{q}} + \mathbf{K}\mathbf{q} = \mathbf{F}(t)
\end{equation}

onde $\mathbf{M}$ é a matriz de massa, $\mathbf{C}$ é a matriz de amortecimento, $\mathbf{G}$ é a matriz giroscópica, $\mathbf{K}$ é a matriz de rigidez, $\mathbf{q}$ é o vetor de coordenadas generalizadas e $\mathbf{F}(t)$ é o vetor de forças externas.

\subsection{Modelagem por Elementos Finitos}

A modelagem por elementos finitos tem se mostrado uma ferramenta eficaz para a análise de sistemas rotativos complexos. Lalanne e Ferraris \cite{lalanne1998} apresentam uma abordagem sistemática para a modelagem de rotores utilizando elementos finitos, incluindo a consideração de efeitos não lineares e acoplamentos entre diferentes componentes.

\section{Análise de Vibrações}

A análise de vibrações constitui a base para o diagnóstico de falhas em máquinas rotativas. Randall \cite{randall2011} apresenta uma abordagem abrangente para a análise de vibrações em aplicações industriais, incluindo técnicas de processamento de sinais e interpretação de espectros.

\subsection{Análise no Domínio da Frequência}

A análise no domínio da frequência é fundamental para o diagnóstico de falhas, permitindo a identificação de componentes espectrais características de diferentes tipos de falhas. A transformada de Fourier é a ferramenta mais utilizada para esta análise:

\begin{equation}
X(f) = \int_{-\infty}^{\infty} x(t) e^{-j2\pi ft} dt
\end{equation}

\subsection{Limitações da Análise Linear}

A análise espectral convencional baseada na transformada de Fourier apresenta limitações quando aplicada a sinais não lineares ou não estacionários. Estas limitações motivaram o desenvolvimento de técnicas mais avançadas, como os espectros de ordem superior.

\section{Espectros de Ordem Superior}

Os espectros de ordem superior (HOS) constituem uma extensão da análise espectral convencional, sendo capazes de capturar informações sobre não linearidades e acoplamentos de fase que não são visíveis nos espectros de segunda ordem.

\subsection{Fundamentos Teóricos}

Segundo Nikias e Petropulu \cite{nikias1993}, os espectros de ordem superior são definidos como momentos estatísticos de ordem superior de sinais estocásticos. O bi-espectro, que é o espectro de terceira ordem, é definido como:

\begin{equation}
B(f_1, f_2) = E[X(f_1)X(f_2)X^*(f_1 + f_2)]
\end{equation}

onde $X(f)$ é a transformada de Fourier do sinal $x(t)$ e $E[\cdot]$ denota o valor esperado.

\subsection{Aplicações em Diagnóstico de Falhas}

Fackrell et al. \cite{fackrell1995} demonstraram a eficácia dos bi-espectros para a análise de sinais de vibração, mostrando que estas técnicas são capazes de identificar não linearidades características de diferentes tipos de falhas em sistemas rotativos.

\section{Diagnóstico de Falhas em Máquinas Rotativas}

O diagnóstico de falhas em máquinas rotativas tem evoluído significativamente com o desenvolvimento de técnicas avançadas de análise de sinais e aprendizado de máquina.

\subsection{Tipos de Falhas}

As falhas mais comuns em máquinas rotativas incluem:

\begin{itemize}
    \item \textbf{Trincas:} Podem ocorrer devido a fadiga, corrosão ou sobrecarga, alterando significativamente a rigidez do rotor.
    \item \textbf{Desalinhamentos:} Resultam em forças desbalanceadas e podem causar vibrações excessivas.
    \item \textbf{Desbalanceamento:} Causa vibrações síncronas com a velocidade de rotação.
    \item \textbf{Falhas em rolamentos:} Podem resultar em vibrações com frequências características.
\end{itemize}

\subsection{Modelagem de Falhas}

A modelagem de falhas em sistemas rotativos requer a consideração de efeitos não lineares. Wensing \cite{wensing1997} apresenta uma revisão abrangente sobre a dinâmica de rotores com trincas, incluindo modelos para o comportamento de "respiração" das trincas.

\section{Aprendizado de Máquina Aplicado ao Diagnóstico}

O uso de técnicas de aprendizado de máquina tem se mostrado promissor para a automação do processo de diagnóstico de falhas. Lei et al. \cite{lei2020} apresentam uma revisão sistemática das aplicações de aprendizado de máquina em prognóstico de saúde de máquinas.

\subsection{Classificadores Supervisionados}

Widodo e Yang \cite{widodo2007} demonstram a eficácia de Support Vector Machines (SVM) para o diagnóstico de falhas, enquanto Kankar et al. \cite{kankar2011} aplicam diferentes algoritmos de aprendizado de máquina para o diagnóstico de falhas em rolamentos.

\section{Lacunas na Literatura}

Apesar dos avanços significativos na área, algumas lacunas permanecem:

\begin{itemize}
    \item Falta de metodologias integradas que combinem simulação numérica, análise HOS e aprendizado de máquina.
    \item Necessidade de validação de técnicas HOS com dados simulados de alta fidelidade.
    \item Desenvolvimento de frameworks automatizados para diagnóstico de falhas.
\end{itemize}

\section{Contribuições Propostas}

Este trabalho propõe contribuir para o preenchimento dessas lacunas através do desenvolvimento de uma metodologia integrada que combine:

\begin{enumerate}
    \item Simulação numérica de alta fidelidade utilizando a biblioteca ROSS.
    \item Análise de espectros de ordem superior para extração de características.
    \item Validação física através de comparação com literatura.
    \item Aplicação de aprendizado de máquina para classificação automática.
\end{enumerate}


\label{sec:intro} % A label to reference this section later

Hello, world! This is my first document created with \LaTeX. 
You write text just like this, and LaTeX handles all the formatting. 
It automatically indents new paragraphs.

You can easily add text formatting like \textbf{bold text} or \textit{italic text}.

\section{Why LaTeX is Useful} % Creates a second section

One of the main strengths of LaTeX is typesetting mathematics. You can write simple inline equations like Einstein's famous $E = mc^2$, or more complex displayed equations like Euler's identity:
\[
e^{i\pi} + 1 = 0
\]

It's also great for creating structured lists:
\begin{itemize}
    \item Item one is for organization.
    \item Item two is for beautiful math.
    \item Item three is for automated references. As we saw in section~\ref{sec:intro}, cross-referencing is simple!
\end{itemize}

This is just the beginning of what is possible.

\end{document}